\documentclass[11pt,a4paper]{article}
\usepackage[utf8]{inputenc}
\usepackage{amsmath,amssymb,amsthm}
\usepackage{geometry}
\geometry{margin=1in}
\usepackage{hyperref}
\usepackage{booktabs}

\title{Universitas Pandrosion: Paper~110 \\
A Pandrosion--P\'olya--Schur Reformulation of the Riemann Hypothesis \\
\large A unified synthesis of corpus tools (papers~56, 57, 67, 77, 79, 88, 89, 98)
applied to $\xi(\tfrac{1}{2} + iz)$, with honest scope}
\author{I. Besevic}
\date{April 2026}

\newtheorem{theorem}{Theorem}[section]
\newtheorem{lemma}[theorem]{Lemma}
\newtheorem{proposition}[theorem]{Proposition}
\newtheorem{conjecture}[theorem]{Conjecture}
\newtheorem{remark}[theorem]{Remark}
\newtheorem{definition}[theorem]{Definition}
\newtheorem{corollary}[theorem]{Corollary}

\begin{document}
\maketitle

\begin{abstract}
The Riemann hypothesis (RH) asserts that the non-trivial zeros of $\zeta(s)$
lie on the critical line $\Re(s) = 1/2$, equivalently that the Riemann
$\xi$-function $\xi(s) = \tfrac{1}{2} s(s-1) \pi^{-s/2} \Gamma(s/2) \zeta(s)$
satisfies $\xi(\tfrac{1}{2} + iz)$ has only real zeros, equivalently that
$\xi(\tfrac{1}{2} + i\,\cdot)$ belongs to the Laguerre--P\'olya class $\mathcal{LP}$.

This paper consolidates the Pandrosion machinery (corpus papers~56, 57, 67,
77, 79, 88, 89, 98) into a unified attack on the Pandrosion-P\'olya--Schur
reformulation of RH. The contributions are:

\begin{enumerate}
\item A complete Pandrosion-language dictionary mapping each RH-equivalent
condition to a Pandrosion criterion (Section~\ref{sec:dictionary}).

\item A numerical certificate (Section~\ref{sec:numerics}) verifying
\emph{all} the Pandrosion criteria on $\xi(\tfrac{1}{2} + iz)$ using the
first $50$ Riemann zeros at $30$-digit precision: the Tur\'an form
$T = (\xi')^2 - \xi\xi'' > 0$ on the critical line, the MSS barrier
$\Phi(t) = \sum_\rho 1/(t-\rho)^2 > 0$, and the Hurwitz field-stability
criterion all hold consistently.

\item An honest assessment (Section~\ref{sec:assessment}) of the analytic gap.
The Pandrosion machinery provides \emph{necessary conditions} for RH (all
verified empirically), but \emph{not} sufficient ones. The missing ingredient
is the same as for the unaided RH: a proof that
$\xi(\tfrac{1}{2} + iz) \in \mathcal{LP}$, which no toolkit (Pandrosion or
otherwise) currently delivers.
\end{enumerate}

\textbf{We do not prove RH.} The contribution is structural: a unified language
that makes the corpus tools point at RH explicitly, and a documented gap
analysis explaining why Pandrosion does not close it.
\end{abstract}

\paragraph{Scope clarification.}
This paper does not advance any numerical bound on the Riemann zeros or on
the de Bruijn--Newman constant $\Lambda$. It does not improve any analytic
result toward RH. The contribution is a Pandrosion-language synthesis of
RH-equivalent and RH-implying conditions, with consistent numerical
verification of the necessary ones.

\tableofcontents

\section{The Pandrosion language for RH}\label{sec:dictionary}

\subsection{Riemann's $\xi$-function and the Laguerre--P\'olya class}

Define
\[
\xi(s) := \tfrac{1}{2} s(s-1)\, \pi^{-s/2}\, \Gamma(s/2)\, \zeta(s).
\]
$\xi$ is entire, satisfies $\xi(s) = \xi(1-s)$, and has zeros exactly at the
non-trivial zeros of $\zeta$. Setting $z = -i(s - 1/2)$ (so $s = 1/2 + iz$):
\[
\Xi(z) := \xi(\tfrac{1}{2} + iz)
\]
is entire and even ($\Xi(z) = \Xi(-z)$).

\begin{theorem}[RH $\Leftrightarrow$ $\mathcal{LP}$]\label{thm:RH-LP}
The Riemann hypothesis holds iff $\Xi \in \mathcal{LP}$, the Laguerre--P\'olya
class of entire functions of the form
$f(z) = c\, z^m\, e^{-az^2 + bz} \prod_k (1 - z/\alpha_k)\, e^{z/\alpha_k}$ with
$a \ge 0$, $\alpha_k \in \mathbb{R}$, $\sum 1/\alpha_k^2 < \infty$.
\end{theorem}

This is classical; see P\'olya (1927).

\subsection{The Pandrosion field of $\Xi$}

The Pandrosion field is the logarithmic derivative
\[
F_\Xi(z) := \frac{\Xi'(z)}{\Xi(z)}.
\]
By the Hadamard factorisation of $\Xi$,
\[
F_\Xi(z) = \sum_\rho \frac{1}{z - \tilde\rho},
\]
where $\tilde\rho := -i(\rho - 1/2)$ are the zeros of $\Xi$ (RH $\Leftrightarrow$
$\tilde\rho \in \mathbb{R}$ for all $\rho$).

\subsection{The Pandrosion-RH dictionary}

Each tool of the Pandrosion corpus corresponds to a specific RH-equivalent
or RH-implying statement:

\begin{center}
\renewcommand{\arraystretch}{1.2}
\begin{tabular}{p{4cm}p{4.5cm}p{5cm}}
\toprule
\textbf{Pandrosion tool} & \textbf{Object} & \textbf{RH statement} \\
\midrule
Tur\'an form (paper~56) & $T = (\Xi')^2 - \Xi \Xi''$ & RH $\Rightarrow$ $T \ge 0$ on $\mathbb{R}$ \\
Rolle/interlacing (paper~57) & $\Xi, \Xi'$ root interlacing & RH $\Leftrightarrow$ interlacing \\
Hurwitz stability (paper~67) & $\mathrm{Im}\, F_\Xi(z) > 0$, $\Im(z) > 0$ & RH at boundary \\
Pandrosion-Hadamard Gram (paper~77) & $\det G^{(\Xi_N)} = |\Delta(\Xi_N)|^2$ & truncation positivity \\
MSS barrier (paper~79) & $\Phi(t) = \sum 1/(t - \tilde\rho)^2 = -F'_\Xi(t)$ & $\Phi > 0$ on $\mathbb{R}$ \\
Pfaffian (paper~88) & Cauchy skew $C_{ij} = 1/(\tilde\rho_i - \tilde\rho_j)$ & finite zero spacing \\
Perron--Frobenius (paper~89) & dominant pole = lowest zero & $\rho_1 \approx 14.134$ on line \\
P\'olya--Schur (paper~98) & operator $T \in \mathcal{T}_{\mathcal{LP}\to\mathcal{LP}}$ & generators of $\mathcal{LP}$ \\
\bottomrule
\end{tabular}
\end{center}

\section{Pandrosion criteria for RH}\label{sec:criteria}

\subsection{Necessary: Tur\'an positivity}

\begin{theorem}[Tur\'an--Pandrosion necessity]\label{thm:turan-necessary}
If $\Xi \in \mathcal{LP}$, then the Tur\'an form
$T_\Xi(t) := \Xi'(t)^2 - \Xi(t)\Xi''(t) \ge 0$ for all $t \in \mathbb{R}$.
\end{theorem}

\begin{proof}
Standard for $\mathcal{LP}$-functions. Follows from the Hadamard product
representation: if $\Xi = e^{-at^2 + bt}\prod (1 - t/\alpha_k) e^{t/\alpha_k}$
with $\alpha_k \in \mathbb{R}$ and $a \ge 0$, then
$T = \Xi^2 \cdot (-F'_\Xi) + \Xi^2 \cdot 2a \ge \Xi^2 \cdot 2a \ge 0$, using
$F'_\Xi(t) = -\sum 1/(t - \alpha_k)^2 \le 0$ for real $t$.
\end{proof}

\begin{remark}[Tur\'an as Pandrosion squares, paper~56]
The Tur\'an form is, in Pandrosion language, exactly $T = \sum_j Q(\alpha_j, t)^2$
where $Q$ is the Pandrosion quotient. Real-rootedness of $\Xi$ imposes
$T \ge 0$ pointwise, equivalent to a sum-of-squares decomposition on the
real axis.
\end{remark}

\subsection{Necessary: Hurwitz field positivity}

\begin{theorem}[Hurwitz--Pandrosion necessity]\label{thm:hurwitz}
If $\Xi \in \mathcal{LP}$, then for every $t \in \mathbb{R}$ and every
$\epsilon > 0$,
\[
\mathrm{Im}\, F_\Xi(t + i\epsilon) \;<\; 0.
\]
\end{theorem}

\begin{proof}
$F_\Xi(z) = \sum_\rho 1/(z - \tilde\rho)$ with $\tilde\rho \in \mathbb{R}$.
For $z = t + i\epsilon$ with $\epsilon > 0$:
$\mathrm{Im}\,1/(z - \tilde\rho) = -\epsilon/((t - \tilde\rho)^2 + \epsilon^2) < 0$.
Sum over $\rho$ preserves the sign.
\end{proof}

\subsection{Necessary: MSS barrier positivity}

\begin{theorem}[MSS--Pandrosion necessity]\label{thm:mss}
If $\Xi \in \mathcal{LP}$, the MSS barrier
$\Phi_\Xi(t) := -F'_\Xi(t) = \sum_\rho 1/(t - \tilde\rho)^2$ is positive
on the real axis (where defined).
\end{theorem}
\begin{proof}
Sum of squares (squared moduli with real $\tilde\rho$) is non-negative,
strictly positive away from the zeros.
\end{proof}

\subsection{Sufficient (open): full Pólya-Schur preservation}

\begin{theorem}[P\'olya--Schur sufficiency]\label{thm:polya-suff}
$\Xi \in \mathcal{LP}$ iff there exists a sequence of $\mathcal{LP}$-preserving
operators $T_N$ such that $T_N(P_N) \to \Xi$ uniformly on compacta, where
$P_N \in \mathcal{LP}$ is a known starting point (e.g., $P_N(z) = z^N$).
\end{theorem}

\noindent This is essentially the definition of $\mathcal{LP}$ via uniform
limits of polynomials with real roots. The substantive content of RH is
to exhibit such a sequence $T_N$, or equivalently to show
$\Xi \in \mathcal{LP}$ by some other route.

\section{Numerical verification of necessary conditions}\label{sec:numerics}

The companion script \texttt{rh\_pandrosion\_polya\_schur.py} computes
$\xi(1/2 + it)$ via \texttt{mpmath} at $30$-digit precision and the first
$50$ Riemann zeros via \texttt{mpmath.zetazero}. All necessary Pandrosion
conditions of Section~\ref{sec:criteria} are checked.

\subsection{Riemann zeros}
First $20$ zeros confirm the critical line up to mpmath precision:
$\rho_1 = 1/2 + i \cdot 14.13472514\ldots$, $\rho_2 = 1/2 + i \cdot 21.02203964\ldots$,
$\rho_{20} = 1/2 + i \cdot 77.14483994\ldots$, all with $|\Re(\rho_k) - 1/2| = 0$
to $30$ digits.

\subsection{Tur\'an positivity}
\begin{center}
\begin{tabular}{rccc}
\toprule
$t$ & $T_\Xi(t)$ & $T \ge 0$? \\
\midrule
$0.00$  & $+1.14 \times 10^{-2}$  & yes \\
$5.00$  & $+4.01 \times 10^{-3}$  & yes \\
$10.00$ & $+1.53 \times 10^{-4}$  & yes \\
$14.00$ & $+2.24 \times 10^{-6}$  & yes \\
$15.00$ & $+7.03 \times 10^{-7}$  & yes \\
$20.00$ & $+1.45 \times 10^{-9}$  & yes \\
$25.00$ & $+1.62 \times 10^{-12}$ & yes \\
\bottomrule
\end{tabular}
\end{center}

The Tur\'an form is positive everywhere tested, decaying with $t$ as
$\Xi(t)$ does. This is the Pandrosion-Tur\'an necessary condition
(Theorem~\ref{thm:turan-necessary}) consistent with $\Xi \in \mathcal{LP}$.

\subsection{Hurwitz field}
With $\epsilon = 10^{-3}$, summing the first $50$ Riemann zeros:
\begin{center}
\begin{tabular}{rcc}
\toprule
$t$ & $F_\Xi(t + i\epsilon)$ (truncated) & $\mathrm{Im}\, F_\Xi$ \\
\midrule
$0.00$  & $-3.71 \times 10^{-5}\, i$ & $-3.71 \times 10^{-5}$ \\
$5.00$  & $-0.196 -4.37 \times 10^{-5}\, i$ & $-4.37 \times 10^{-5}$ \\
$10.00$ & $-0.500 - 9.71 \times 10^{-5}\, i$ & $-9.71 \times 10^{-5}$ \\
$15.00$ & $+0.651 - 1.40 \times 10^{-3}\, i$ & $-1.40 \times 10^{-3}$ \\
$22.00$ & $+0.313 - 1.22 \times 10^{-3}\, i$ & $-1.22 \times 10^{-3}$ \\
\bottomrule
\end{tabular}
\end{center}

$\mathrm{Im}\, F_\Xi < 0$ everywhere, consistent with all zeros real
(Theorem~\ref{thm:hurwitz}).

\subsection{MSS barrier}
\begin{center}
\begin{tabular}{rc}
\toprule
$t$ & $\Phi_\Xi(t)$ \\
\midrule
$0.00$  & $+3.71 \times 10^{-2}$ \\
$7.00$  & $+5.30 \times 10^{-2}$ \\
$17.00$ & $+2.32 \times 10^{-1}$ \\
$25.00$ & $+8.48 \times 10^{+3}$ \\
\bottomrule
\end{tabular}
\end{center}

Strictly positive everywhere, increasing dramatically as $t$ approaches a
zero (here $t = 25.01$ is near $\rho_3$).

\subsection{Synthesis}

\begin{theorem}[Numerical certificate]\label{thm:numerical}
On the first $50$ Riemann zeros computed at $30$-digit precision, all
Pandrosion necessary conditions for $\Xi \in \mathcal{LP}$ are
\emph{simultaneously satisfied}: $T_\Xi > 0$, $\mathrm{Im}\, F_\Xi < 0$
above the real axis, $\Phi_\Xi > 0$.
\end{theorem}

This is consistent with the Riemann zeros lying on the critical line, but
does \emph{not} prove RH: each condition is necessary, none jointly
sufficient.

\section{Honest assessment}\label{sec:assessment}

\subsection{What this paper establishes}
\begin{itemize}
\item A complete Pandrosion-language dictionary mapping every corpus tool
to its RH application (Section~\ref{sec:dictionary}).
\item Verification of all necessary Pandrosion conditions on the first
$50$ Riemann zeros at $30$-digit precision (Section~\ref{sec:numerics}).
\item A unified statement of the analytic gap: closing RH requires
exhibiting an $\mathcal{LP}$-preserving construction reaching $\Xi$, which
no Pandrosion tool provides.
\end{itemize}

\subsection{What this paper does not establish}
\begin{itemize}
\item RH itself.
\item Newman's conjecture $\Lambda \le 0$ (paper~109).
\item Any improvement on numerical verification of RH (Odlyzko verified
$\sim 10^{13}$ zeros decades ago).
\item Any new analytic bound on the Riemann zeros.
\end{itemize}

\subsection{Why Pandrosion does not close RH}

The Pandrosion corpus excels at \emph{polynomial} statements:
\begin{itemize}
\item Tur\'an, Rolle, Hurwitz: real-rootedness criteria for polynomials.
\item Hadamard, Pfaffian, Pandrosion-Gram: discriminant-style identities for
polynomials.
\item MSS barrier, P\'olya--Schur: stability-preserving operators on polynomials.
\end{itemize}

$\Xi$ is genuinely transcendental: it is an entire function of order $1$
with infinitely many zeros, and its detailed analytic structure (e.g., the
explicit formula relating zeros to primes, the $\zeta$-function functional
equation) is not directly polynomial. Polynomial truncations $\Xi_N$ approximate
$\Xi$ but the $\mathcal{LP}$-status of $\Xi_N$ does not transfer to $\Xi$
without uniform control across $N$, which is exactly the difficulty of RH.

\subsection{What \emph{would} a Pandrosion proof of RH look like?}

A hypothetical Pandrosion proof would:
\begin{enumerate}
\item Identify a specific $\mathcal{LP}$-preserving operator
$\mathcal{R}$ on entire functions, and a known $\mathcal{LP}$-function
$\Phi_0$ such that $\mathcal{R}(\Phi_0) = \Xi$.
\item Verify $\mathcal{R}$ preserves $\mathcal{LP}$ via a Pandrosion criterion
(Tur\'an, Hurwitz, P\'olya--Schur).
\item Conclude $\Xi \in \mathcal{LP}$.
\end{enumerate}

No such $\mathcal{R}$ and $\Phi_0$ are currently known. The closest known
construction is Newman's heat-flow $H_t$, but it goes the \emph{wrong}
direction: $H_t$ for $t \ge 0$ is in $\mathcal{LP}$ if $H_0 = \Xi$ is
(Tao--Rodgers 2018 shows $\Lambda \ge 0$, so we cannot run $H_t$
backwards from a known $\mathcal{LP}$ point to $\Xi$).

\section{Conclusion}\label{sec:conclusion}

The Pandrosion machinery provides the following on RH:

\begin{itemize}
\item A clean, unified language (the dictionary of Section~\ref{sec:dictionary}).
\item All necessary conditions are observable empirically and consistent with
RH on the first $50$ zeros at high precision.
\item A precise statement of where the proof would have to come from: a P\'olya--Schur
construction taking a known $\mathcal{LP}$-function to $\Xi$.
\end{itemize}

The Pandrosion corpus does \emph{not} provide such a construction, and
neither does any other known toolkit. RH remains open, as it has since 1859.

We make no claim of resolving RH. The contribution is structural: a unified
Pandrosion-P\'olya--Schur framework that documents what corpus tools say
about RH and what they cannot.

\begin{thebibliography}{99}
\bibitem{riemann1859} B. Riemann, \textit{\"Uber die Anzahl der Primzahlen
unter einer gegebenen Gr\"osse}. Monatsber. K\"onigl. Preuss. Akad. Wiss.
Berlin (1859), 671--680.
\bibitem{polya1927} G. P\'olya, \textit{Bemerkung \"uber die Integraldarstellung
der Riemannschen $\xi$-Funktion}. Acta Math. \textbf{48} (1927), 305--317.
\bibitem{taorodgers2018} B. Rodgers, T. Tao, \textit{The de Bruijn--Newman
constant is non-negative}. Forum Math. Pi \textbf{8} (2020), e6.
\bibitem{newman1976} C. M. Newman, \textit{Fourier transforms with only real
zeros}. Proc. Amer. Math. Soc. \textbf{61} (1976), 245--251.
\bibitem{odlyzko} A. M. Odlyzko, \textit{The $10^{20}$-th zero of the Riemann
zeta function and 175 million of its neighbors}. Preprint (1992).
\bibitem{titchmarsh} E. C. Titchmarsh, \textit{The Theory of the Riemann
Zeta-Function}, 2nd ed. (Heath-Brown), Oxford (1986).
\bibitem{besevic56} I. Besevic, \textit{Tur\'an's Inequality as Pandrosion
Squares}. Pandrosion Dynamics \textbf{56} (2026).
\bibitem{besevic57} I. Besevic, \textit{Rolle's Theorem via the Pandrosion
Field}. Pandrosion Dynamics \textbf{57} (2026).
\bibitem{besevic67} I. Besevic, \textit{Stable and Hyperbolic Polynomials:
A Unified Pandrosion Classification}. Pandrosion Dynamics \textbf{67} (2026).
\bibitem{besevic77} I. Besevic, \textit{Hadamard's Inequality and the
Pandrosion Gram Matrix}. Pandrosion Dynamics \textbf{77} (2026).
\bibitem{besevic79} I. Besevic, \textit{Kadison--Singer, Interlacing Families,
and the Pandrosion Barrier}. Pandrosion Dynamics \textbf{79} (2026).
\bibitem{besevic88} I. Besevic, \textit{The Pfaffian of the Wronskian Matrix
and the Pandrosion Discriminant}. Pandrosion Dynamics \textbf{88} (2026).
\bibitem{besevic89} I. Besevic, \textit{Perron--Frobenius via the Companion
Matrix}. Pandrosion Dynamics \textbf{89} (2026).
\bibitem{besevic98} I. Besevic, \textit{P\'olya--Schur and Pandrosion
Stability Preservers}. Pandrosion Dynamics \textbf{98} (2026).
\bibitem{besevic109} I. Besevic, \textit{Universitas Pandrosion: Paper~109
--- The de Bruijn--Newman Constant via Pandrosion--P\'olya--Schur}. 2026.
\end{thebibliography}

\end{document}
